% Text tokenization and embeddings, shared by the NNDS CNN deck and the NN
% transformer deck. Moved from Lecture_CNNs.tex by Claude (2026-11-16).

% Text pipeline rewritten step by step at Simone's request (2026-10-30), with
% new figures in shared/assets/images/text/{text_pipeline,tokenization,vocabulary,
% embedding_lookup}.tex replacing the old tokenization and embedding figures.
\begin{whiteframe}{Handling text data}

    Text must be preprocessed to be handled by a neural network. We turn it
    into a matrix in three steps:

    \vspace{0.3em}
    \begin{center}
        % Editable text-to-matrix pipeline: tokenize, look up IDs, embed. Each token
% becomes one row of X. Uses the same text and IDs as the next-token slides.
% Written by Claude (2026-10-30), replacing the old embedding figure.
\begingroup
\begin{tikzpicture}[x=1cm,y=1cm,
        tok/.style={draw=drawgray,fill=drawgray_l,rounded corners=1.5pt,line width=0.7pt,
            minimum width=1.35cm,minimum height=0.55cm,inner sep=2pt,font=\small\ttfamily},
        tid/.style={draw=draworange,fill=draworange_l,rounded corners=1.5pt,line width=0.7pt,
            minimum width=0.7cm,minimum height=0.55cm,inner sep=0pt,font=\small},
        ecell/.style={draw=drawblue,fill=drawblue_l,line width=0.5pt,minimum width=0.42cm,
            minimum height=0.55cm,inner sep=0pt},
        edge/.style={-{Stealth[length=2mm]},line width=0.8pt,draw=drawgray},
        steplab/.style={font=\small\bfseries,anchor=south,align=center},
        under/.style={font=\small,anchor=north}]
    \node[draw=drawgray,rounded corners=2pt,line width=0.7pt,inner sep=5pt,
        font=\small\ttfamily] (text) at (0,0) {Check this out};
    \foreach \w/\id [count=\t from 0] in {Check/2,this/16,out/3}{
        \pgfmathsetmacro{\y}{0.75-0.75*\t}
        \node[tok] (w\t) at (3.6,\y) {\w};
        \node[tid] (z\t) at (6.0,\y) {$\id$};
        \foreach \d in {0,...,4}{
            \node[ecell] at (8.3+0.42*\d,\y) {};
        }
        \draw[edge] (w\t.east) -- (z\t.west);
        \draw[edge] (z\t.east) -- (8.0,\y);
        \draw[drawgray,line width=0.7pt] (text.east) -- (2.15,0) |- (w\t.west);
    }
    \draw[edge] (2.15,0.75) -- (w0.west);
    \draw[edge] (2.15,0) -- (w1.west);
    \draw[edge] (2.15,-0.75) -- (w2.west);
    \draw[drawblue,line width=1pt] (8.09,1.025) rectangle (10.19,-1.025);
    \node[steplab] at (2.2,1.25) {1. tokenize};
    \node[steplab] at (4.8,1.25) {2. look up IDs};
    \node[steplab] at (7.35,1.25) {3. embed};
    \node[under] at (0,-1.25) {text};
    \node[under] at (3.6,-1.25) {$T$ tokens};
    \node[under] at (6.0,-1.25) {$\vc{z}\sim(T)$};
    \node[under] at (9.14,-1.25) {$\vc{X}\sim(T,D)$};
\end{tikzpicture}
\endgroup

    \end{center}

    \vspace{0.5em}
    {\raggedright
        By the end of this procedure we have a matrix $\vc{X}\sim(T,D)$ of $T$
        tokens, each represented by $D$ features\NNDSorNN{}{ (we write $\vc{X}\sim(T,D)$ for a matrix with $T$ rows and $D$ columns)}.\par}

\end{whiteframe}

\begin{whiteframe}{Step 1: choosing the tokens}
    \setlength{\parskip}{4pt}

    We first choose how to split (\myalert{tokenize}) the text, e.g., by
    character, by byte, by \myalert{subword} or by word:

    \begin{center}
        % Editable tokenization comparison on one sentence: characters, subwords and
% words, with the resulting sequence length. Subwords are the actual output of
% the GPT-4 tokenizer (tiktoken, cl100k_base). Written by Claude (2026-10-30),
% replacing the old tokenization figure.
\begingroup
\begin{tikzpicture}[x=1cm,y=1cm,
        tk/.style={draw=#1,fill=#1_l,rounded corners=1.5pt,line width=0.6pt,
            minimum height=0.5cm,inner sep=2.5pt,font=\small\ttfamily,anchor=west},
        rowlabel/.style={font=\small\bfseries,anchor=east},
        len/.style={font=\small,anchor=west,text=drawgray}]
    % Visible space drawn with rules: the monospace font has no U+2423 glyph.
    \def\sp{{\color{drawgray}\rule[-0.3ex]{0.5pt}{0.7ex}\rule[-0.3ex]{0.42em}{0.5pt}\rule[-0.3ex]{0.5pt}{0.7ex}}}
    \node[font=\small\ttfamily,anchor=west] at (0,0.95) {Classify this text!};
    \node[rowlabel] at (-0.2,0.95) {Text};
    % Characters (spaces are tokens too).
    \node[rowlabel] at (-0.2,0) {Characters};
    \foreach \ch [count=\n from 0] in {C,l,a,s,s,i,f,y,\sp,t,h,i,s,\sp,t,e,x,t,!}{
        \node[tk=drawblue,minimum width=0.4cm,inner sep=0pt] at (0.42*\n,0) {\ch};
    }
    \node[len] at (8.3,0) {$T=19$};
    % Subwords.
    \node[rowlabel] at (-0.2,-0.8) {Subwords};
    \node[tk=draworange] (s0) at (0,-0.8) {Class};
    \node[tk=draworange] (s1) at ([xshift=2pt]s0.east) {ify};
    \node[tk=draworange] (s2) at ([xshift=2pt]s1.east) {\sp this};
    \node[tk=draworange] (s3) at ([xshift=2pt]s2.east) {\sp text};
    \node[tk=draworange] (s4) at ([xshift=2pt]s3.east) {!};
    \node[len] at (8.3,-0.8) {$T=5$};
    % Words.
    \node[rowlabel] at (-0.2,-1.6) {Words};
    \node[tk=drawgreen] (v0) at (0,-1.6) {Classify};
    \node[tk=drawgreen] (v1) at ([xshift=2pt]v0.east) {this};
    \node[tk=drawgreen] (v2) at ([xshift=2pt]v1.east) {text};
    \node[tk=drawgreen] (v3) at ([xshift=2pt]v2.east) {!};
    \node[len] at (8.3,-1.6) {$T=4$};
\end{tikzpicture}
\endgroup

    \end{center}

    Larger \myalert{tokens} give shorter sequences, but we need many more of
    them to cover all possible texts. Byte tokenization works like character
    tokenization on the $256$ byte values of UTF-8 text.

    \SlideSource{Subwords from the GPT-4 tokenizer (\texttt{cl100k\_base}), via
        \href{https://github.com/openai/tiktoken}{tiktoken}.}

\end{whiteframe}

\begin{frame}{Step 1: building the vocabulary}
    \setlength{\parskip}{4pt}

    The set of tokens we can handle is the \myalert{vocabulary} or
    \myalert{dictionary}. Its size $V$ depends on the type of token:

    \medskip
    \begin{columns}[T,onlytextwidth]
        \begin{column}{0.32\textwidth}
            \begin{tcolorbox}[nndsblue,courseheight={text-vocab}{}]
                \textbf{Characters or bytes}\\
                A small fixed set, e.g., the $256$ byte values.
            \end{tcolorbox}
        \end{column}
        \begin{column}{0.32\textwidth}
            \begin{tcolorbox}[nndsred,courseheight={text-vocab}{}]
                \textbf{Words}\\
                The $V$ most frequent words of a training corpus, plus a
                special \textit{unknown} token.
            \end{tcolorbox}
        \end{column}
        \begin{column}{0.32\textwidth}
            \begin{tcolorbox}[nndsgreen,courseheight={text-vocab}{}]
                \textbf{Subwords}\\
                Reusable pieces learned from a corpus, e.g., with
                \myalert{byte pair encoding} (BPE).
            \end{tcolorbox}
        \end{column}
    \end{columns}

    \medskip
    Larger tokens need a larger vocabulary: a few hundred characters, but
    tens of thousands of words or subwords.

\end{frame}

\begin{frame}{Subword embeddings}
    \setlength{\parskip}{3pt}
    \setlength{\abovedisplayskip}{6pt}
    \setlength{\belowdisplayskip}{6pt}

    Word-level vocabularies can miss unseen words. Character tokenization
    produces longer sequences. With \myalert{subword tokenization}, we reuse
    pieces across words, for example (illustrative):
    \[
        \begin{array}{rcl}
            \texttt{playing} & \longmapsto & (\texttt{play},\texttt{ing}) \\
            \texttt{played}  & \longmapsto & (\texttt{play},\texttt{ed})
        \end{array}
    \]

    Each piece has an ID and a row in $\vc{E}\sim(V,D)$.
    Both words share the embedding for \texttt{play}.

    Unseen words can split into known pieces, and frequent words can stay whole.
    Learned pieces need not match linguistic word parts.

    \SlideSource{Sennrich, Haddow and Birch (2016),
        \href{https://aclanthology.org/P16-1162/}{\textit{Neural Machine Translation of Rare Words with Subword Units}}, ACL.}

\end{frame}

\begin{frame}{Byte pair encoding}
    \setlength{\parskip}{3pt}
    \setlength{\abovedisplayskip}{5pt}
    \setlength{\belowdisplayskip}{5pt}

    We can learn the vocabulary with \myalert{byte pair encoding} (BPE):
    \begin{enumerate}
        \item Start the vocabulary with the single characters.
        \item Merge the most frequent adjacent pair into a new token, \myalert{added to the vocabulary}.
        \item Repeat until the vocabulary reaches the chosen size $V$.
    \end{enumerate}

    Toy corpus: \texttt{hug} $\times5$, \texttt{hugs} $\times3$,
    \texttt{pug} $\times2$. We count pairs within words:
    \[
        \begin{array}{rcll}
            (\texttt{u},\texttt{g})  & \longmapsto & \texttt{ug}  & \quad 10\text{ occurrences} \\
            (\texttt{h},\texttt{ug}) & \longmapsto & \texttt{hug} & \quad 8\text{ occurrences}
        \end{array}
    \]
    The vocabulary grows from $\{\texttt{g},\texttt{h},\texttt{p},\texttt{s},\texttt{u}\}$
    to $\{\ldots,\texttt{ug}\}$ and then $\{\ldots,\texttt{ug},\texttt{hug}\}$.

    On new text, we apply the fixed merge rules in their learned order.
    \myalert{Byte-level BPE} starts from $256$ byte values, so it can represent
    any UTF-8 text.

    \SlideSource{\href{https://huggingface.co/learn/llm-course/en/chapter6/5}{Hugging Face: Byte-Pair Encoding tokenization}.}

\end{frame}

% Figure: shared/assets/images/text/vocabulary.tex (Claude, 2026-10-30).
\begin{whiteframe}{Step 2: from tokens to IDs}
    \setlength{\parskip}{4pt}

    Each token in the text is converted into a discrete ID
    $z_t\in\{0,\ldots,V-1\}$, its position in the vocabulary, giving us an
    integer sequence $\vc{z}\sim(T)$:

    \begin{center}
        % Editable word-level vocabulary: a table of IDs and tokens, and the lookup
% of two texts, one with a typo that falls back to the unknown token. IDs
% match the embedding-lookup slide (Check = 2, this = 16, out = 3).
% Written by Claude (2026-10-30).
\begingroup
\begin{tikzpicture}[x=1cm,y=1cm,
        vrow/.style={minimum height=0.42cm,inner sep=0pt,font=\small},
        tok/.style={draw=drawgray,fill=drawgray_l,rounded corners=1.5pt,line width=0.6pt,
            minimum height=0.5cm,inner sep=2.5pt,font=\small\ttfamily},
        tid/.style={draw=draworange,fill=draworange_l,rounded corners=1.5pt,line width=0.6pt,
            minimum width=0.7cm,minimum height=0.5cm,inner sep=0pt,font=\small},
        edge/.style={-{Stealth[length=2mm]},line width=0.8pt,draw=drawgray},
        head/.style={font=\small\bfseries,anchor=south}]
    % Vocabulary table.
    \node[head] at (0.9,1.45) {Vocabulary};
    \foreach \id/\tk [count=\n from 0] in {
            0/{\textless unk\textgreater},1/{the},2/{Check},3/{out},
            \vdots/{\vdots},16/{this},\vdots/{\vdots},V{-}1/{zebra}}{
        \pgfmathsetmacro{\y}{1.2-0.42*\n}
        \node[vrow,anchor=east,text=draworange] at (0.45,\y-0.21) {$\id$};
        \node[vrow,anchor=west,font=\small\ttfamily] at (0.7,\y-0.21) {\tk};
    }
    \draw[drawgray,line width=0.6pt,rounded corners=2pt] (-0.25,1.25) rectangle (2.1,-2.21);
    \draw[drawgray!60,line width=0.4pt] (0.58,1.25) -- (0.58,-2.21);

    % Two lookups.
    \foreach \row/\yy in {0/0.75,1/-1.05}{
        \ifcase\row
            \node[tok,anchor=west] (a\row) at (3.4,\yy) {Check};
            \node[tok,anchor=west] (b\row) at ([xshift=2pt]a\row.east) {this};
            \node[tok,anchor=west] (c\row) at ([xshift=2pt]b\row.east) {out};
        \or
            \node[tok,anchor=west] (a\row) at (3.4,\yy) {Check};
            \node[tok,anchor=west] (b\row) at ([xshift=2pt]a\row.east) {this};
            \node[tok,anchor=west,draw=drawred,text=drawred] (c\row) at ([xshift=2pt]b\row.east) {ouut};
        \fi
    }
    \foreach \row/\yy/\ids in {0/0.75/{2,16,3},1/-1.05/{2,16,0}}{
        \foreach \i [count=\n from 0] in \ids{
            \node[tid] (i\row\n) at (7.9+0.8*\n,\yy) {$\i$};
        }
        \draw[edge] ([xshift=0.15cm]c\row.east) -- ([xshift=-0.1cm]i\row0.west);
    }
    \node[font=\footnotesize,text=drawgray,anchor=north] at (5.4,-1.45)
        {\texttt{ouut} is not in the vocabulary};
\end{tikzpicture}
\endgroup

    \end{center}

    A word outside the vocabulary, such as a typo, becomes the unknown token.

\end{whiteframe}

\begin{frame}{Step 3: dense embeddings}
    \setlength{\parskip}{4pt}
    \setlength{\abovedisplayskip}{7pt}
    \setlength{\belowdisplayskip}{7pt}

    For embedding we learn an \myalert{embedding matrix} $\vc{E}\sim(V,D)$, with one row for each of the $V$ tokens in the vocabulary. Embedding element $z_t$ of the sequence is just a row lookup:
    \[
        \idx{\vc{X}}{t,:}=\idx{\vc{E}}{z_t,:}
    \]

    We initialize the embeddings randomly and train them with the network, or take them from some pre-trained model or algorithm.
    For example, GloVe and Word2Vec were two classic algorithms for this purpose.

    Pre-trained embeddings also inherit the biases of their training text:
    for example, word analogies can reproduce gender stereotypes about
    professions.

    \SlideSource{Bolukbasi et al. (2016), \href{https://arxiv.org/abs/1607.06520}{\textit{Man is to Computer Programmer as Woman is to Homemaker? Debiasing Word Embeddings}}, NeurIPS.}

\end{frame}

% Figure: shared/assets/images/text/embedding_lookup.tex (Claude, 2026-10-30).
\begin{whiteframe}{Visualizing the embedding procedure}

    \vspace{0.5em}
    \begin{center}
        % Editable embedding lookup: the IDs of "Check this out" select rows of
% E ~ (V,D), which become the rows of X ~ (T,D), the input of a 1D CNN (NNDS) or a transformer (NN).
% Written by Claude (2026-10-30), replacing the old embedding figure.
\begingroup
\begin{tikzpicture}[x=1cm,y=1cm,
        ecell/.style={draw=drawgray!70,line width=0.4pt,minimum width=0.4cm,
            minimum height=0.36cm,inner sep=0pt},
        tid/.style={draw=draworange,fill=draworange_l,rounded corners=1.5pt,line width=0.6pt,
            minimum width=0.7cm,minimum height=0.5cm,inner sep=0pt,font=\small},
        edge/.style={-{Stealth[length=2mm]},line width=0.8pt,draw=drawgray},
        lab/.style={font=\small},
        note/.style={font=\footnotesize,text=drawgray}]
    % Embedding matrix: displayed rows, with the ID of each row on the left.
    \foreach \id/\fill [count=\n from 0] in {
            0/white,1/white,2/drawblue_l,3/drawviolet_l,4/white,
            dots/white,16/draworange_l,dots/white,V{-}1/white}{
        \pgfmathsetmacro{\y}{1.6-0.36*\n}
        \def\tmp{dots}
        \ifx\id\tmp
            \node[note] at (1.0,\y) {$\vdots$};
        \else
            \foreach \d in {0,...,4}{
                \node[ecell,fill=\fill] at (0.2+0.4*\d,\y) {};
            }
            \node[note,anchor=east] at (-0.1,\y) {$\id$};
        \fi
    }
    \node[lab,anchor=south] at (1.0,1.9) {$\vc{E}\sim(V,D)$};

    % X: one row per token.
    \begin{scope}[shift={(4.2,0)}]
        \foreach \fill/\w/\id [count=\t from 0] in {drawblue_l/Check/2,draworange_l/this/16,drawviolet_l/out/3}{
            \foreach \d in {0,...,4}{
                \node[ecell,fill=\fill] at (0.2+0.4*\d,0.76-0.36*\t) {};
            }
            \node[note,anchor=west] at (2.1,0.76-0.36*\t) {\texttt{\w} ($z_{\t}=\id$)};
        }
        \node[lab,anchor=south] at (1.0,1.9) {$\vc{X}\sim(T,D)$};
    \end{scope}
    % Row copies.
    \draw[edge,draw=drawblue] (2.1,0.88) to[out=0,in=180] (4.1,0.76);
    \draw[edge,draw=draworange] (2.1,-0.56) to[out=0,in=180] (4.1,0.40);
    \draw[edge,draw=drawviolet] (2.1,0.52) to[out=0,in=180] (4.1,0.04);
    % Next layer.
    \node[lab,draw=drawblue,fill=drawblue_l!25!white,rounded corners=2pt,
        inner sep=5pt] (cnn) at (10.7,0.4) {\NNDSorNN{1D CNN}{Transformer}};
    \draw[drawblue,line width=1pt] (4.2,0.94) rectangle (6.2,-0.14);
    \draw[edge] (8.75,0.4) -- (cnn.west);
\end{tikzpicture}
\endgroup

    \end{center}

    \vspace{0.6em}
    For $\vc{z}=(2,16,3)$, each ID selects one row of $\vc{E}$. With custom embeddings, we train
    $\vc{E}$ together with the rest of the network\NNDSorNN{, e.g., a 1D CNN for text
    classification}{}.

\end{whiteframe}

\begin{frame}{Sequence padding}
    \setlength{\parskip}{4pt}
    \setlength{\abovedisplayskip}{7pt}
    \setlength{\belowdisplayskip}{7pt}

    Two texts may have different lengths. Their embedded matrices
    $\vc{X}_1\sim(T_1,D)$ and $\vc{X}_2\sim(T_2,D)$ share the same feature dimension $D$.

    For $T_1>T_2$, we can build a padded mini-batch:
    \[
        X=\operatorname{stack}\!\left(
        \vc{X}_1,\begin{bmatrix}\vc{X}_2\\\vc{0}\end{bmatrix}
        \right),
        \qquad X\sim(2,T_1,D).
    \]

    We pad to the maximum sequence length in the mini-batch.
    The added zero rows are placeholders.

    We keep a \myalert{validity mask} of real positions, so padded positions
    do not contribute to pooling or to the loss.

\end{frame}

